% Copyright (c) 2022-2026 MKM Research Labs. Licensed under the MIT License; see LICENSE.
\documentclass[11pt]{article}
\usepackage[utf8]{inputenc}
\usepackage{graphicx}\usepackage{longtable}\usepackage{array}\usepackage{booktabs}
\usepackage{xcolor}\usepackage{parskip}\usepackage{enumitem}\setlist{nosep}
\usepackage{float}\usepackage{lastpage}\usepackage{fancyhdr}\usepackage{hyperref}
\usepackage{amsmath}\usepackage{amsfonts}\usepackage{amssymb}
\usepackage{calc}
\usepackage[toc,page]{appendix}
\usepackage{subcaption}
\usepackage{geometry}\geometry{a4paper, total={170mm,247mm}, left=20mm, top=25mm, bottom=25mm}
\usepackage{setspace}\onehalfspacing
\setcounter{tocdepth}{3}\setcounter{secnumdepth}{3}
% --- MKM Research Labs house style (mirrors src/governance/latex.py) ---
\definecolor{mkmblue}{HTML}{1565C0}
\hypersetup{colorlinks=true,linkcolor=mkmblue,citecolor=mkmblue,urlcolor=mkmblue}
\pagestyle{fancy}\fancyhf{}
\fancyhead[L]{\small\textcolor{gray}{MKM Research Labs}}
\fancyhead[R]{\small\textcolor{gray}{Interest Rate Swaption --- v1.0}}
\fancyfoot[C]{\small\textcolor{gray}{CONFIDENTIAL --- PROPRIETARY}}
\fancyfoot[R]{\small\textcolor{gray}{Page \thepage\ of \pageref{LastPage}}}
\renewcommand{\headrulewidth}{0.4pt}\renewcommand{\footrulewidth}{0.4pt}

\begin{document}

% --- Title page (mirrors latex.py title_page + legal_notice) ---
\begin{titlepage}\begin{center}
\vspace*{1cm}
\IfFileExists{../Miscellaneous/logo.jpg}{\includegraphics[width=0.4\textwidth]{../Miscellaneous/logo.jpg}\\[1.5cm]}{\vspace*{1cm}}
{\Huge\bfseries\textcolor{mkmblue}{Interest Rate Swaption}\par}\vspace{0.5cm}
{\Large\textcolor{gray}{Pricing and Risk Methodology}\par}\vspace{2cm}
{\large MKM Research Labs\par}\vspace{0.5cm}
{\large Version 1.0 --- \today\par}\vspace{0.5cm}
{\large MKM-ModelRisk\par}
\vspace{2cm}\textcolor{red}{\textbf{CONFIDENTIAL --- PROPRIETARY}}
\end{center}
\vfill
\begin{minipage}{\textwidth}\footnotesize\raggedright
\textbf{Legal notice}\par\medskip
PROPRIETARY AND CONFIDENTIAL\par\medskip
This document, including all algorithms, models, methodology, formulas, calculations, and intellectual concepts contained herein, constitutes the exclusive intellectual property and confidential trade secrets of MKM Research Labs (``MKM'').\par\medskip
Any usage, reproduction, distribution, modification, or disclosure of this document or any portion thereof, without the express prior written authorisation from MKM Research Labs is strictly prohibited and constitutes an infringement of intellectual property rights.\par\medskip
The document is provided on an ``as is'' basis. MKM Research Labs makes no representations or warranties of any kind, express or implied, regarding its accuracy, reliability, or suitability for any particular purpose.\par\medskip
All rights reserved. \textcopyright{} 2021--2026 MKM Research Labs.\par\medskip
Governed by the laws of the United Kingdom.
\end{minipage}
\end{titlepage}

\tableofcontents
\clearpage

\section{Purpose and Scope}
This document provides an outline of how to price and risk an Interest Rate Swaption. \par

\subsection{Support}
For any assistance, problems or suggestions regarding this methodology, contact MKM Research Labs:
\begin{itemize}
  \item Website: \href{https://example.com}{example.com}
  \item Email: \href{mailto:support@example.com}{support@example.com}
\end{itemize}

\section{Background}

Swap options, or swaptions, are options on interest rate swaps and are another popular type of interest rate option. They give the holder the right but not the obligation to enter into an interest rate swap at a specified time in the future. \par

Corporates use swaptions to hedge their issuance of debt, with a guarantee that the fixed rate of interest they will pay will not exceed some level. Instead of paying upfront, a corporate can enter a forward-starting swap that locks the corporate into a swap agreement. Swaptions follow other options strategies, such as hedging foreign exchange exposure with an FX option. \par

Consider a swaption where the holder has the right to pay a rate $s_K$ and receive float on a swap of principle $L$ that will last $n$ years (swap tenor) starting in $T$ years (option maturity) with $m$ payments and notional $L$. \par

Swaptions come in three types, notably European (option exercises at maturity), American (option exercises anytime prior to maturity) and Bermudian (exercisable at specific dates prior to maturity, usually annually).  The assumption at this juncture is that all swaptions are European, which is by far the most liquid market. \par


One aspect of swaption trading is how the market-makers quote in two dimensions of option maturity on the y-axis and swap length, or tenor, on the x-axis.  Different users operate in different sectors of the swaption grid.  For example pension funds purchase short-term options to hedge their long-term liabilities.  In trading parlance, pension funds are operating in the top-right quadrant.  Structured products structured on the back of long-term infrastructural  projects such as bridges might operate in the bottom-right quadrant where they buy long-dated options into long tenor swaps.  At the peak of the derivative madness in early 2000s, it was possible to see 60 year options exercising into 60 year swaps.  Normal trading occurs on the top left to bottom right diagonal where liquidity can concentrate in the 5s5s that is a five-year option into a five-year swap. \par

Unlike equities that build a curve across strikes, swaption quotes are at-the-money only, though when building an implied volatility, introducing a skew will be essential.  The only problem is where to source OTM option volatility quotes. \par

The valuation of a swaption that is covered extensively by Hull and White\cite{art_hull_white_2003} \par

\section{Valuation of a Swaption}
We now start to push the Swaption forward into a Black-Scholes model for pricing the models.  Consider we have a call option on a Swaption with unit currency notional.  The purchaser has the right (but not the obligation) to buy protection on a Swaption that starts at time $T>0$, maturing at $T^*>T$ for a spread K.  We have the following:- \par

\begin{table}[ht]
	\centering
	\begin{tabular}{ll}
		Parameter & Description\\
		\toprule[1pt]
		$N$ Notional of contract \\[2ex]
		$F$ & Forward swap rate of a n-year swap starting at time T \\ [2ex]
		$K$ & swaption strike\\ [2ex]
		$r$ risk free rate to T \\[2ex]
		$c_0$ & Present Value of the Call Swaption \\ [2ex]
		$p_0$ & Present Value of the Call Swaption \\ [2ex]
		$\sigma_{bs}$ & standard deviation of $ln F$ \\ [2ex]
		$\sigma_{b}$ & standard deviation of $ F$ \\ [2ex] 
		\bottomrule[1pt]
	\end{tabular}
	\caption{Model inputs for valuing the Swaption option }
	\label{tab:model_inputs_Swaption}
\end{table}

Depending on whether the distribution for the forward swap rate is Lognormal or Normal depends on its prevailing rate and volatility.  If the rate is unlikely to go negative and has space to "move without needing to go negative", the Black Scholes is fine.  If the rate is close to zero and has a tendancy or opportunity to go negative, then Bachelier is better.\par

\subsection{Black Scholes formula}
\begin{align}
	c_0 & =  N e^{-rT}\left [ F N(d_1) - K N(d_2) \right ] \\
	p_0 & =  N e^{-rT}\left [ K N(-d_2) - F N(-d_1) \right ] \\
	d_1& = \displaystyle \frac{Ln(F/K)+(\sigma_{bs}^2/2)T}{\sigma_{bs} \sqrt{T}} \\
	d_2 &= d_1 - \sigma \sqrt{T}
\end{align}

\subsection{Bachelier formula}

\begin{align}
	c_0 & =  N e^{-rT}\left [ (F-K) N(d_1) + \sigma_{b} \sqrt{T}N(d_1) \right ] \\
	p_0 & =  N e^{-rT}\left [ (K-F) N(d_1) + \sigma_{b} \sqrt{T}N(d_1) \right ] \\
	d_1& = \displaystyle \frac{(F-K)}{\sigma \sqrt{T}} \\
\end{align}
\section{Introduction to implied volatility surface}

Unlike other asset classes that operate in two dimensions, a swaption volatility surface is a three-dimensional plot across strike, option maturity and swap tenor.  A crucial property of the implied volatility surface is the absence of arbitrage. \par
All volatility surfaces must comply with the arbitrage-free conditions such as the following \par

\begin{itemize}
	\item Static arbitrage free condition: Static arbitrage free condition makes it impossible to invest nothing today and receive positive return tomorrow. 
	\item Calendar arbitrage free condition: The cost of a calendar spread should be positive.
	\item Vertical (spread) arbitrage free condition: The cost of a vertical spread across strikes should be positive.
	\item Horizontal (butterfly) arbitrage free condition: The cost of a butterfly spread across strikes should be positive
\end{itemize}

Note, there is no calendar arbitrage in swaption volatility surfaces as swaptions with different (maturities x tenors) have different underlying swaps and are associated with other underlying indices.  We need to treat them independently, though maturity combinations such as (6x6), (7x5) and (5x7) are closely aligned; just not connected rigidly courtesy of any arbitrage rule. \par


\section{Creating Implied Volatility}

There are a couple of techniques used to build a full volatility surface.  The Ricardo Rebonato book\cite{book_rebonato_2009} provides the outline of SABR that is largely adopted by the industry where there are limited number of quotes (usually five but possibily can drop to three.  For larger sets, the complexity of the SABR model we found was an overkill given the amount of data provided by ICE and a more straightforward approach of cubic spline  \par

\subsection{SABR approach}

Hagan in 2002\cite{art_hagan_2002} introduced singular perturbation techniques to derive the implied volatility for a plain-vanilla (European) swaption.  Within the context of Black Scholes implied volatility:- \par

\begin{align}
	T &= \text{Option Maturity} \\
	K &= \text{Option Strike} \\
	f &= \text{Forward rate for swap starting at T of term t} \\
	\sigma(K, T) &= \Xi(K) \dfrac{z}{x(z)} \Psi(K)	
\end{align}

where \par
\begin{align}
	\Xi(K)  &= \dfrac{\alpha}{\left(f K\right)^{\dfrac{1-\beta}{2}} \left[ 1 + \dfrac{(1-\beta)^2}{24} ln^2 \dfrac{f}{K} + \dfrac{(1-\beta)^4}{1920} ln^4\dfrac{f}{K}\right] }\\
	\Psi(K) &= 1 + \left[ \dfrac{(1-\beta)^2\alpha^2}{24(f)K)^{1-\beta}} + \dfrac{\rho\beta\nu\alpha}{4 (f K)^{\frac{1-\beta}{2}}} + \nu^2 \dfrac{2-3\rho^2}{24} \right].T\\
	z       &= \dfrac{\nu}{\alpha} \left(f K\right)^{\dfrac{1-\beta}{2}} ln\left(\dfrac{f}{K} \right)\\
	\chi(z)    &= ln \left[ \dfrac{\sqrt{1-2\rho z + z^2} + z-\rho}{1-\rho} \right] \\
		\sigma(K, T) &\approx \Xi(K) \dfrac{z}{x(z)} \Psi(K)	
\end{align}

	The four model parameters $(\alpha, \beta , \rho, \nu )$ above have the following intuitive meaning:- \par

\begin{itemize}
	\item  $\alpha$ \quad The volatility that is tied closely to the at-the-money volatility value
	\item  $\beta$ \quad Defines the nature of how the forward rate varies (lognormal is one and normal is zero)
	\item  $\rho$ \quad  Describes the “skew” or average slope of the volatility curve across strikes
	\item  $\nu$ \quad Describes the “smile” or curvature/convexity of the volatility curve across the strike for a given term and tenor
\end{itemize}

For ATM forward options where the distribution is somewhere between normal and log-normal

\begin{equation}
	\sigma(K=f) = \displaystyle \dfrac{\alpha}{f^{1-\beta}}. \times \left ( 1 + \left [\dfrac{ \alpha^2 (1 - \beta)^2}{24f^{2(1-\beta)}} + \frac{\rho\beta\nu\alpha}{4f^{(1-\beta)}} +\frac{(2-3 \rho^2)\nu^2}{24} \right ].T \right)
\end{equation}

We have solved for $\alpha$ and set $\beta$ as a static model input.  This leaves us to find the optimized value of 
$(\rho, \nu)$ by minimizing the distance between the SABR model output volatilities and market volatilities across all strikes for each term and tenor.  Industry norm tends to use the Levenberg-Marquardt least-squares optimization routine. \par
\begin{equation}
	min \sum_{i=1}^{n} \left (  \sigma_i^{SABR} - \sigma_{i}^{Market} \right )^2
\end{equation}

The SABR calibration process will then need to be repeated for each option maturity and swap tenor \par

\subsection{Calibrating ICE Market Data}

We found that the amount of ICE swaption market data so comprehensive across maturities, strikes and tenors that the SABR model could not converge.  \par

\begin{table}[ht]
	\centering
	\begin{tabular}{c|c|c|c|c}
	Tenor & Maturity & Strike & NormalVol & LogVol \\ [1ex]
	5Y & 1Y & -0.49888 & 134.82 & 0 \\ [1ex]
	5Y & 1Y & -0.24888 & 133.03 & 0 \\ [1ex]
	5Y & 1Y & 0.00112 & 131.5 & 0 \\ [1ex]
	5Y & 1Y & 0.25112 & 130.32 & 173.48 \\ [1ex]
	5Y & 1Y & 0.50112 & 129.54 & 127.61 \\ [1ex]
	5Y & 1Y & 0.75112 & 129.24 & 106.05 \\ [1ex]
	5Y & 1Y & 1.00112 & 129.49 & 92.9 \\ [1ex]
	5Y & 1Y & 1.25112 & 130.31 & 84.01 \\ [1ex]
	5Y & 1Y & 1.50112 & 131.74 & 77.65 \\ [1ex]
	5Y & 1Y & 1.75112 & 133.75 & 72.98 \\ [1ex]
	5Y & 1Y & 2.00112 & 136.31 & 69.49 \\ [1ex]
	5Y & 1Y & 2.25112 & 139.37 & 66.85 \\ [1ex]
	5Y & 1Y & 2.50112 & 142.85 & 64.84 \\ [1ex]
	5Y & 1Y & 2.75112 & 146.69 & 63.29 \\ [1ex]
	5Y & 1Y & 3.00112 & 150.81 & 62.1 \\ [1ex]
	5Y & 1Y & 3.25112 & 155.17 & 61.18 \\ [1ex]
	5Y & 1Y & 3.50112 & 159.71 & 60.46 \\ [1ex]
	5Y & 1Y & 3.75112 & 164.38 & 59.9 \\ [1ex]
	5Y & 1Y & 4.00112 & 169.17 & 59.47 \\ [1ex]
	5Y & 1Y & 4.25112 & 174.03 & 59.13 \\ [1ex]
	5Y & 1Y & 4.50112 & 178.96 & 58.86 \\ [1ex]
	\end{tabular}
	\label{tab:ice_data}
\end{table}

\subsection{Cubic Spline}

Spline fitting (or interpolation) provides a method to draw a smooth curve through $n+1$ points $(x_0,y_0), \dots ,(x_n,y_n))$ and we look for polynomials $p_0, \dots ,p_{n-1}$ where $p_i(x)$ operates within the interval $[x_i,x_{i+1}]$.  Once we have defined the polynomials, we define $f(x): f(x_i)=p(x_i)$ across all integers \par

\begin{align}
	p_i(x) &= a_i(x-x_i)^3 + b_i(x-x_i)^2 + c_i(x-x_i)+d_i \\
	\text{Where for } i &= 0, \dots , n-1\\
	a_ i &= \displaystyle \frac{b_{i+1}-b_i}{3h_i} \\
	c_i &= \displaystyle \frac{k_i}{h_i} - \displaystyle \frac{h_i}{3}(b_{i+1} + 2b_i) \\
	d_i &= y_i 
\end{align}

Define $h_i = x_{i+1} - x_i$ and $k_i = y_{i+1} - y_i$.  The $b_i$ coefficient is defined by a matrix operate as follows \par

\begin{align}
	\text{Define for } i &= 1, \dots ,n \\
	h_ i &= \displaystyle \frac{h_i}{h_{i-1} + h_i} \\
	r_ i &= \displaystyle \frac{h_{i-1}}{h_{i-1} + h_i} = 1 - s_i \\
	v_i &= \displaystyle \frac {3 \left (\frac{k_i}{h_i} - \frac{k_{i-1}}{h_{i-1}} \right ) }{h_{i-1}+h_i}
\end{align}

We define the matrix $U=[u_{ij}]$ to be a $n+1 \times n+1$ matrix 
\begin{equation}
	u_{ij} = \begin{cases}
		r_i  &  i = j+1 \\
		2 & i = j \\
		s_i & i = j-1 \\
		0  & \text{otherwise}
	\end{cases}
\end{equation}

Define $V=[v_i]$ that is a $n+1$ vector using the formula above. \par
Define $B=[bi]$ that is defined as \par
\begin{equation}
	B=U^{-1} V
\end{equation} 


\clearpage

\section{Document History}
\begin{longtable}{llll}
\caption{Document history}\\
\toprule
\textbf{Date} & \textbf{Version} & \textbf{Description} & \textbf{Author} \\
\midrule
\endfirsthead
\toprule
\textbf{Date} & \textbf{Version} & \textbf{Description} & \textbf{Author} \\
\midrule
\endhead
31-Oct-2022 & v1.0 & Internal beta release (library v1.0 Beta) & MKM-ModelRisk \\
\bottomrule
\end{longtable}

\clearpage

\bibliographystyle{plain}
\bibliography{../Miscellaneous/bibliography}

\clearpage
\begin{appendices}

This page intentionally left blank

\end{appendices}
\end{document}
